% Generated from locale/tr/content/incompleteness/representability-in-q/beta-function.tex; do not hand-edit.


\olfileid[tr]{inc}{req}{bet}
\olsection{Beta Fonksiyonu Lemması}


Toplama, çarpma ve $\Char{=}$ elimizde olduğunda ilkel özyinelemeyi
gerçekleştirebildiğimizi göstermek için dizileri işleyen fonksiyonlar
geliştirmemiz gerekir. (Üs alma da elimizde olsaydı işimiz daha kolay
olurdu.) İlkel özyinelemeye sahipken ``$n$'inci asal'' gibi şeyleri
tanımlayabilir ve oldukça doğrudan bir kodlama seçebilirdik. Fakat burada
ilkel özyinelememiz yoktur---gerçekte ilkel özyinelemeyi minimizasyon
kullanarak yapabildiğimizi göstermek istiyoruz---bu yüzden daha yaratıcı
olmalıyız.

\begin{lem}
\ollabel{lem:beta}
Her $a_0$, \dots,~$a_n$ dizisi için, her $i\le n$'de
$\beta(d,i)=a_i$ olacak bir $d$ sayısının bulunduğu bir $\beta(d,i)$
fonksiyonu vardır. Üstelik $\beta$, yalnızca bileşke ve düzenli minimizasyon
kullanılarak temel fonksiyonlardan tanımlanabilir.
\end{lem}

$d$'nin $\tuple{a_0,\dots,a_n}$ dizisini kodladığını ve $\beta(d,i)$'nin
$i$'inci elemanı döndürdüğünü düşünün. (Bu ``kodlama''nın daha önce
öğrendiğimiz asal kuvvetleri kodlamasını \emph{kullanmadığına} dikkat edin!)
Lemma oldukça sınırlıdır; dizileri art arda ekleyebildiğimizi, eleman
ekleyebildiğimizi, hatta bileşke ve düzenli minimizasyonla tanımlanabilen
fonksiyonları kullanarak $a_0$, \dots,~$a_n$'den $d$'yi
\emph{hesaplayabildiğimizi} bile söylemez. Yalnızca her dizinin
``kodlanacağı'' bir ``kod çözme'' fonksiyonunun bulunduğunu söyler.

$\beta$ gösterimi G\"odel'e aittir. Yineleyelim: Lemmayı ispatlamanın zor
bölümü, görünüşte sınırlı olan kaynaklarla, yani yalnızca bileşke ve
minimizasyonla uygun bir $\beta$ tanımlamaktır; ancak toplama, çarpma ve
$\Char{=}$ kullanmamıza izin vardır. Bu lemmayı ispatlamanın çeşitli yolları
bulunur; en temizlerinden biri hâlâ G\"odel'in özgün yöntemidir. Bu yöntem,
Sunzi Teoremi (geleneksel adıyla ``Çin Kalan Teoremi'') denen sayı-kuramsal
bir olguyu kullanır.

\begin{defn}
İki $a$ ve $b$ doğal sayısı, en büyük ortak bölenleri $1$ ise ancak ve ancak
\emph{aralarında asaldır}; başka bir deyişle başka hiçbir ortak bölenleri
yoktur.
\end{defn}

\begin{defn}
$a$ ve $b$ doğal sayıları, $c\mid(a-b)$ ise ancak ve ancak
\emph{$c$ modülüne göre denktir}; bunu $a\equiv b\mod c$ ile yazarız. Yani
$a$ ile $b$, $c$'ye bölündüğünde aynı kalanı verir.
\end{defn}

Sunzi Teoremi şöyledir:
\begin{thm}
$x_0$, \dots,~$x_n$'in (çifter çifter) aralarında asal olduğunu varsayın.
$y_0$, \dots,~$y_n$ herhangi sayılar olsun. O zaman aşağıdakileri sağlayan
bir $z$ sayısı vardır:
\begin{align*}
z & \equiv y_0 \mod x_0 \\
z & \equiv y_1 \mod x_1 \\
& \vdots  \\
z & \equiv y_n \mod x_n.
\end{align*}
\end{thm}

Sunzi Teoremi'ni şöyle kullanacağız: $x_0$, \dots,~$x_n$ sırasıyla
$y_0$, \dots,~$y_n$'den büyükse, $z$'yi $\tuple{y_0,\dots,y_n}$ dizisinin
kodu alabiliriz. $y_i$'yi geri elde etmek için yalnızca $z$'yi $x_i$'ye
bölüp kalanı almamız gerekir. Bu kodlamayı kullanmak için $x_0$, \dots,
$x_n$'e uygun değerler bulmalıyız.

Birkaç gözlem bu konuda bize yardımcı olacaktır. $y_0$, \dots,~$y_n$
verildiğinde
\begin{align*}
j &= \max(n, y_0 + 1, \dots, y_n + 1), \\
m &= \lcm(1,\dots,j),
\end{align*}
olsun ve
\begin{align*}
x_0 & = 1 + m \\
x_1 & = 1 + 2 \cdot m \\
x_2 & = 1 + 3 \cdot m \\
& \vdots  \\
x_n & = 1 + (n+1) \cdot m
\end{align*}
olsun. O zaman iki şey doğrudur:
\begin{enumerate}
\item\ollabel{rel-prime} $x_0,\dots,x_n$ aralarında asaldır.
\item\ollabel{less} Her $i$ için $y_i<x_i$ olur.
\end{enumerate}
\olref{rel-prime}'ın doğru olduğunu görmek için $p$ bir asal sayı olup
$p\mid x_i$ ve $p\mid x_k$ ise $p\mid 1+(i+1)m$ ve
$p\mid 1+(k+1)m$ olduğuna dikkat edin. Fakat bu durumda $p$ bunların
farkını böler:
\[
(1 + (i+1)m) - (1+ (k+1)m) = (i-k) m.
\]
$p$, $1+(i+1)m$'yi böldüğünden aynı zamanda $m$'yi bölemez (bölseydi ilk
bölmede kalan $1$ olurdu). Dolayısıyla $p$, $(i-k)m$'yi böldüğü için
$p$ aynı zamanda $i-k$'yi böler. Fakat $\left|i-k\right|$ en çok $n$'dir ve $j\geq n$
seçtik; öyleyse bu $p\mid m$ sonucunu verir ve yine çelişkiye ulaşırız.
Demek ki hem $x_i$'yi hem de $x_k$'yi bölen hiçbir asal sayı yoktur.
\olref{less} maddesi kolaydır: $y_i<j\leq m<x_i$ olur.

Şimdi $\beta$ fonksiyonu lemmasını ispatlayalım. $0$, ardıl, toplama,
çarpma, $\Char{=}$, izdüşümler ve bunlardan bileşke ile düzenli fonksiyonlara
uygulanan minimizasyon kullanılarak tanımlanan her fonksiyonu
kullanabileceğimizi anımsayın. Karakteristik fonksiyonu böyle
tanımlanabiliyorsa bir bağıntıyı da kullanabiliriz. Daha önce olduğu gibi bu
bağıntıların Boole birleşimleri ve sınırlı niceleme altında kapalı olduğunu
gösterebiliriz; örneğin:
\begin{align*}
\fn{not}(x) & \defis \Char{=}(x,0)\\
\bmin{x \leq z}{R(x,y)} & \defis \umin{x}{(R(x,y) \lor x = z)}\\
\bexists{x \leq z}{R(x,y)} & \defiff R(\bmin{x \leq z}{R(x,y)}, y)
\end{align*}
Ardından aşağıdakilerin hepsinin de ilkel özyineleme olmadan
tanımlanabildiğini gösterebiliriz:
\begin{enumerate}
\item Çiftleme fonksiyonu, $J(x,y) = \frac{1}{2}[(x+y)(x+y+1)] + x$;
% maybe explain more what is going on here, a bit confusing.
\item izdüşüm fonksiyonları
\begin{align*}
K(z) & = \bmin{x \leq z}{\bexists{y \leq z}{z = J(x,y)}},\\
L(z) & = \bmin{y \leq z}{\bexists{x \leq z}{z = J(x,y)}};
\end{align*}
\item küçüktür bağıntısı $x<y$;
\item bölünebilirlik bağıntısı $x\mid y$;
% \item $x \tsub y$
% \item $\fn{Prime}(x)$
% \item Assuming $p$ is prime, the relation ``$x$ is a power of $p$'':
% \[
% \bforall{y \leq x}{(y \mid x \lif y = 1 \lor y = x)}.
% \]
\item $y$, $x$'e bölündüğünde kalanı döndüren $\fn{rem}(x,y)$ fonksiyonu.
\end{enumerate}
Şimdi
\begin{align*}
\beta^*(d_0,d_1,i) & = \fn{rem}(1+(i+1) d_1,d_0) \text{ ve}\\
\beta(d,i) & = \beta^*(K(d),L(d),i).
\end{align*}
olarak tanımlayın. İstediğimiz fonksiyon budur. Yukarıdaki gibi
$a_0,\dots,a_n$ verildiğinde
\[
j = \max(n,a_0+1,\dots,a_n+1),
\]
olsun ve $d_1=\lcm(1,\dots,j)$ alınsın. Yukarıdaki \olref{rel-prime}
uyarınca $1+d_1$, $1+2d_1$, \dots, $1+(n+1)d_1$ sayılarının aralarında asal
olduğunu; \olref{less} uyarınca bunların sırasıyla $a_0,\dots,a_n$'den
büyük olduğunu biliyoruz. Sunzi Teoremi uyarınca her $i$ için
\[
d_0 \equiv a_i \mod (1+(i+1)d_1)
\]
olacak bir $d_0$ değeri vardır; dolayısıyla ($d_1$, $a_i$'den büyük olduğu
için)
\[
a_i = \fn{rem}(1+(i+1)d_1,d_0).
\]
$d=J(d_0,d_1)$ olsun. O zaman her $i\le n$ için
\begin{align*}
\beta(d,i) & =  \beta^*(d_0,d_1,i) \\
& =  \fn{rem}(1+(i+1) d_1,d_0) \\
& =  a_i
\end{align*}
olur; istediğimiz buydu. Böylece $\beta$-fonksiyonu lemmasının ispatı
tamamlanır.

\begin{prob}
  $x<y$, $x\mid y$ bağıntıları ile $\fn{rem}(x,y)$ fonksiyonunun ilkel
  özyineleme olmadan tanımlanabildiğini gösterin. $0$, ardıl, toplama,
  çarpma, $\Char{=}$, izdüşümler, sınırlı minimizasyon ve niceleme
  kullanabilirsiniz.
\end{prob}

